% =============================================================================
%                    WEEK 2: TEXT REPRESENTATION TECHNIQUES
% =============================================================================
\section{Week 2: Text Representation Techniques}

% -----------------------------------------------------------------------------
%                              2.1 TOPIC SUMMARY
% -----------------------------------------------------------------------------
\subsection{Topic Summary}

\begin{keybox}[Learning Objectives]
By the end of this week, you will be able to:
\begin{enumerate}
  \item Explain the \textbf{distributional hypothesis}: words in similar contexts have similar meanings.
  \item Represent words/documents as vectors using \textbf{term-document} and \textbf{word-word} co-occurrence matrices.
  \item Compare word vectors using \textbf{dot product} and \textbf{cosine similarity}.
  \item Describe weighting methods (\textbf{tf-idf}, \textbf{PMI}) and why they improve on raw counts.
  \item Understand the motivation for \textbf{dense embeddings} (e.g., Word2Vec) over sparse count vectors.
\end{enumerate}
\end{keybox}

\separator

\subsubsection*{The Big Picture}

This week answers: \textit{``How can we represent word meaning computationally?''}

Key insight: instead of treating words as opaque symbols, we represent them as \textbf{vectors} in a high-dimensional space where \textbf{similar meanings} correspond to \textbf{nearby points}.

\begin{itemize}
  \item \textbf{Count-based}: build vectors from co-occurrence statistics (sparse, interpretable).
  \item \textbf{Prediction-based}: learn vectors by training a classifier (dense, generalises better).
\end{itemize}

\separator

\subsubsection*{What Examiners Like to Ask}

\begin{itemize}
  \item State the \textbf{distributional hypothesis} and give an example.
  \item Explain why \textbf{raw dot product} is problematic; derive or apply \textbf{cosine similarity}.
  \item Compute or interpret \textbf{tf-idf} weights (especially why common words get weight $\approx 0$).
  \item Describe \textbf{Skip-Gram with negative sampling}: training objective, role of sigmoid, why scalable.
  \item Discuss \textbf{bias in embeddings} and the \textbf{parallelogram analogy} method (and its caveats).
\end{itemize}

\separator

% -----------------------------------------------------------------------------
%                 2.2 OVERVIEW + CORE CONCEPTS + EXTENDED THEORY
% -----------------------------------------------------------------------------
\subsection{Overview + Core Concepts + Extended Theory}

\subsubsection*{Roadmap}
\begin{enumerate}
  \item Desiderata for word meaning (synonymy, similarity, relatedness, connotation)
  \item The distributional hypothesis
  \item Co-occurrence matrices: term-document vs.\ word-word
  \item Similarity metrics: dot product and cosine
  \item Weighting schemes: tf-idf
  \item Dense embeddings: Word2Vec (Skip-Gram with negative sampling)
  \item Properties, analogies, and bias in embeddings
\end{enumerate}

\bigseparator

% --- Lexical Semantics Desiderata ---
\subsubsection{Lexical Semantics: Desiderata for Word Meaning}

\begin{defbox}[Key Relations Between Word Senses]
\begin{center}
\renewcommand{\arraystretch}{1.3}
\begin{tabular}{@{} l p{9cm} @{}}
\toprule
\textbf{Relation} & \textbf{Description / Example} \\
\midrule
\textbf{Polysemy} & One lemma, multiple senses: \code{mouse} (rodent vs.\ device) \\
\textbf{Synonymy} & Same meaning in some contexts: \code{couch}/\code{sofa}, \code{car}/\code{automobile} \\
\textbf{Antonymy} & Opposites differing on one feature: \code{hot}/\code{cold}, \code{up}/\code{down} \\
\textbf{Similarity} & Share meaning elements but not synonyms: \code{car}/\code{bicycle} \\
\textbf{Relatedness} & Associated via semantic field: \code{coffee}/\code{cup} (related, not similar) \\
\textbf{Connotation} & Affective meaning: \code{copy} (neutral) vs.\ \code{knockoff} (negative) \\
\bottomrule
\end{tabular}
\end{center}
\end{defbox}

\textbf{Connotation dimensions} (Osgood et al., 1957):
\begin{itemize}
  \item \textbf{Valence}: pleasantness (\code{love} = high, \code{toxic} = low)
  \item \textbf{Arousal}: emotional intensity (\code{frenzy} = high, \code{mellow} = low)
  \item \textbf{Dominance}: control exerted (\code{powerful} = high, \code{weak} = low)
\end{itemize}

\begin{tipbox}[Exam Tip: No Perfect Synonyms]
Even near-synonyms differ in register, formality, or connotation (e.g., \code{water} vs.\ \code{H$_2$O}).
\end{tipbox}

\separator

% --- Distributional Hypothesis ---
\subsubsection{The Distributional Hypothesis}

\begin{defbox}[Distributional Hypothesis]
\textbf{``Words that occur in similar contexts tend to have similar meanings.''}

\hfill --- Zellig Harris (1954); Firth (1957)

\separator

\textbf{Implication:} We can infer word meaning from \textit{usage patterns} rather than dictionary definitions.
\end{defbox}

\begin{exbox}[Example: Inferring Meaning from Context]
Suppose you encounter the unknown word \code{ongchoi} in:
\begin{itemize}
  \item ``Ongchoi is delicious saut\'{e}ed with garlic.''
  \item ``Ongchoi leaves with salty sauces.''
\end{itemize}
Compare with known contexts for \code{spinach}, \code{chard}, \code{collard greens}.

\textbf{Conclusion:} \code{ongchoi} is a leafy green vegetable --- inferred purely from shared context words.
\end{exbox}

\separator

% --- Vector Semantics ---
\subsubsection{Vector Semantics: Words as Points in Space}

\begin{defbox}[Core Idea]
Represent each word as a \textbf{vector} in a high-dimensional space.

\begin{itemize}
  \item \textbf{Similar meanings} $\Rightarrow$ \textbf{nearby vectors}.
  \item Enables graded similarity (not just ``same'' or ``different'').
  \item Allows generalisation to \textit{unseen but similar} words.
\end{itemize}
\end{defbox}

\textbf{Why vectors help ML:}
\begin{itemize}
  \item One-hot features require exact word match between train and test.
  \item Embedding features: similar words have similar vectors $\Rightarrow$ classifier generalises.
\end{itemize}

\separator

% --- Co-occurrence Matrices ---
\subsubsection{Co-occurrence Matrices}

\begin{defbox}[Term-Document Matrix]
Rows = vocabulary words; Columns = documents.

Cell $(w, d)$ = count (or weighted count) of word $w$ in document $d$.

\textbf{Use case:} Information retrieval, document similarity.
\end{defbox}

\begin{defbox}[Word-Word (Term-Term) Matrix]
Rows = vocabulary words; Columns = also vocabulary words.

Cell $(w, c)$ = count of context word $c$ within a $\pm k$ window of target word $w$.

\textbf{Use case:} Capturing fine-grained semantic similarity.
\end{defbox}

\begin{center}
\renewcommand{\arraystretch}{1.2}
\begin{tabular}{@{} l c c @{}}
\toprule
 & \textbf{Term-Document} & \textbf{Word-Word} \\
\midrule
Columns & Documents & Context words \\
Window & Whole document & $\pm k$ words (e.g., $\pm 4$) \\
Typical use & IR, topic models & Word similarity \\
\bottomrule
\end{tabular}
\end{center}

\begin{exbox}[Example: Mini Word-Word Matrix]
Corpus windows:
\begin{itemize}
  \item ``\ldots cherry \textbf{pie}, a traditional \ldots''
  \item ``\ldots strawberry rhubarb \textbf{pie}. Apple \textbf{pie} \ldots''
  \item ``\ldots personal \textbf{digital} assistants \ldots''
  \item ``\ldots \textbf{information} available on the internet \ldots''
\end{itemize}

\begin{center}
\renewcommand{\arraystretch}{1.2}
\begin{tabular}{@{} l c c c @{}}
\toprule
 & \code{a} & \code{computer} & \code{pie} \\
\midrule
\code{cherry} & 1 & 0 & 1 \\
\code{strawberry} & 0 & 0 & 2 \\
\code{digital} & 0 & 1 & 0 \\
\code{information} & 1 & 1 & 0 \\
\bottomrule
\end{tabular}
\end{center}

Observation: \code{cherry} and \code{strawberry} share the \code{pie} dimension; \code{digital} and \code{information} share \code{computer}.
\end{exbox}

\begin{tipbox}[Sparsity]
Full word-word matrices are $|V| \times |V|$ (e.g., $50{,}000 \times 50{,}000$). Most entries are zero $\Rightarrow$ \textbf{sparse vectors}. Efficient sparse-matrix algorithms are essential.
\end{tipbox}

\separator

% --- Similarity Metrics ---
\subsubsection{Similarity Metrics: Dot Product and Cosine}

\begin{defbox}[Dot Product]
\[
\vec{v} \cdot \vec{w} = \sum_{i=1}^{N} v_i w_i
\]

High when vectors have large values in the same dimensions.

\textbf{Problem:} Favours \textit{long} vectors $\Rightarrow$ frequent words dominate.
\end{defbox}

\begin{defbox}[Cosine Similarity]
\[
\cos(\vec{v}, \vec{w}) = \frac{\vec{v} \cdot \vec{w}}{|\vec{v}| \, |\vec{w}|} = \frac{\sum_{i} v_i w_i}{\sqrt{\sum_{i} v_i^2} \; \sqrt{\sum_{i} w_i^2}}
\]

\textbf{Properties:}
\begin{itemize}
  \item Normalises by vector length $\Rightarrow$ removes frequency bias.
  \item Focuses on \textit{direction} (pattern of co-occurrences), not magnitude.
  \item For non-negative vectors: ranges from $0$ (orthogonal) to $1$ (identical direction).
\end{itemize}
\end{defbox}

\begin{exbox}[Example: Cosine Calculation]
Given (from Wikipedia):

\begin{center}
\renewcommand{\arraystretch}{1.2}
\begin{tabular}{@{} l c c c @{}}
\toprule
 & \code{pie} & \code{data} & \code{computer} \\
\midrule
\code{cherry} & 442 & 8 & 2 \\
\code{digital} & 5 & 1683 & 1670 \\
\code{information} & 5 & 3982 & 3325 \\
\bottomrule
\end{tabular}
\end{center}

\[
\cos(\text{cherry}, \text{information}) = \frac{442 \cdot 5 + 8 \cdot 3982 + 2 \cdot 3325}{\sqrt{442^2+8^2+2^2}\;\sqrt{5^2+3982^2+3325^2}} \approx 0.017
\]
\[
\cos(\text{digital}, \text{information}) \approx 0.996
\]

\textbf{Conclusion:} \code{information} is much closer to \code{digital} than to \code{cherry}.
\end{exbox}

\separator

% --- tf-idf ---
\subsubsection{Weighting: tf-idf}

Raw counts overweight frequent but uninformative words (e.g., \code{the}, \code{good}).

\begin{defbox}[tf-idf]
\[
\text{tf-idf}(t, d) = \text{tf}_{t,d} \times \text{idf}_t
\]

\textbf{Term frequency (tf):}
\[
\text{tf}_{t,d} = \begin{cases}
1 + \log_{10}(\text{count}(t,d)) & \text{if count} > 0 \\
0 & \text{otherwise}
\end{cases}
\]

\textbf{Inverse document frequency (idf):}
\[
\text{idf}_t = \log_{10}\left(\frac{N}{\text{df}_t}\right)
\]
where $N$ = total documents, $\text{df}_t$ = number of documents containing $t$.
\end{defbox}

\begin{exbox}[Example: idf in Shakespeare's 37 Plays]
\begin{center}
\renewcommand{\arraystretch}{1.2}
\begin{tabular}{@{} l c c l @{}}
\toprule
\textbf{Word} & \textbf{df} & \textbf{idf} & \textbf{Interpretation} \\
\midrule
\code{Romeo} & 1 & 1.57 & Very discriminative (one play) \\
\code{Falstaff} & 4 & 0.97 & Moderately discriminative \\
\code{good} & 37 & 0 & Appears everywhere $\Rightarrow$ no discriminative power \\
\bottomrule
\end{tabular}
\end{center}

\textbf{Key insight:} If $\text{df}_t = N$, then $\text{idf}_t = \log(1) = 0$, so tf-idf $= 0$.
\end{exbox}

\begin{tipbox}[What Counts as a ``Document''?]
For tf-idf, a ``document'' can be a play, article, paragraph, or even a sentence --- whatever granularity suits the task.
\end{tipbox}

\bigseparator

% --- Word2Vec ---
\subsubsection{Dense Embeddings: Word2Vec}

\begin{defbox}[Sparse vs.\ Dense Vectors]
\begin{center}
\renewcommand{\arraystretch}{1.2}
\begin{tabular}{@{} l l l @{}}
\toprule
 & \textbf{Sparse (count-based)} & \textbf{Dense (learned)} \\
\midrule
Dimensionality & $|V|$ (20k--50k) & 50--1000 \\
Most entries & Zero & Non-zero \\
Synonymy & Distinct dimensions & Similar vectors \\
Generalisation & Limited & Better \\
\bottomrule
\end{tabular}
\end{center}
\end{defbox}

\begin{defbox}[Word2Vec: Skip-Gram with Negative Sampling (SGNS)]
\textbf{Idea:} Learn embeddings by training a classifier to predict whether a word $c$ appears near target $w$.

\textbf{Self-supervision:} Positive examples come from actual context; no manual labels needed.
\end{defbox}

\textbf{Training procedure:}
\begin{enumerate}
  \item For each target word $w$ in the corpus, collect context words $c_1, \ldots, c_L$ within window $\pm k$.
  \item These $(w, c_i)$ pairs are \textbf{positive examples}.
  \item Randomly sample $k$ \textbf{negative examples} $(w, c_{\text{neg}})$ from the vocabulary.
  \item Train a logistic classifier to distinguish positive from negative pairs.
  \item The learned weight matrices become the embeddings.
\end{enumerate}

\begin{defbox}[Skip-Gram Probability Model]
Probability that $c$ is a true context word of $w$:
\[
P(+|w, c) = \sigma(c \cdot w) = \frac{1}{1 + \exp(-c \cdot w)}
\]

\textbf{Sigmoid} converts the dot product into a probability.

For all context words (assuming independence):
\[
P(+|w, c_{1:L}) = \prod_{i=1}^{L} \sigma(c_i \cdot w)
\]
\end{defbox}

\begin{defbox}[Skip-Gram Loss Function (Negative Sampling)]
For one target $w$ with positive context $c_{\text{pos}}$ and $k$ negative samples $c_{\text{neg}_1}, \ldots, c_{\text{neg}_k}$:
\[
\mathcal{L} = -\left[ \log \sigma(c_{\text{pos}} \cdot w) + \sum_{i=1}^{k} \log \sigma(-c_{\text{neg}_i} \cdot w) \right]
\]

\textbf{Goal:} Maximise dot product with true context; minimise dot product with noise.
\end{defbox}

\begin{tipbox}[Why Negative Sampling?]
Full softmax requires normalising over all $|V|$ words at each step --- computationally prohibitive. Negative sampling approximates this by contrasting against a small random sample.
\end{tipbox}

\begin{exbox}[Example: Skip-Gram Training Data]
Sentence: ``\ldots lemon, a \textbf{tablespoon of apricot jam, a} pinch \ldots''

Target: \code{apricot}, window $\pm 2$

\textbf{Positive pairs:} (\code{apricot}, \code{tablespoon}), (\code{apricot}, \code{of}), (\code{apricot}, \code{jam}), (\code{apricot}, \code{a})

\textbf{Negative pairs} (sampled): (\code{apricot}, \code{aardvark}), (\code{apricot}, \code{matrix}), \ldots
\end{exbox}

\textbf{Gradient descent update intuition:}
\begin{itemize}
  \item Move $w$ and $c_{\text{pos}}$ closer (increase their dot product).
  \item Move $w$ and each $c_{\text{neg}}$ apart (decrease their dot products).
\end{itemize}

\begin{defbox}[Two Embedding Matrices]
Skip-Gram learns:
\begin{itemize}
  \item $W$ --- target word embeddings
  \item $C$ --- context word embeddings
\end{itemize}

Common practice: combine them as $\vec{w}_i + \vec{c}_i$ for the final representation.
\end{defbox}

\separator

% --- Properties of Embeddings ---
\subsubsection{Properties of Embeddings}

\textbf{Window size effects:}
\begin{center}
\renewcommand{\arraystretch}{1.2}
\begin{tabular}{@{} l p{10cm} @{}}
\toprule
\textbf{Window} & \textbf{Nearest Neighbours} \\
\midrule
Small ($\pm 2$) & Syntactically/taxonomically similar (e.g., \code{Hogwarts} $\to$ other fictional schools) \\
Large ($\pm 5$) & Topically related (e.g., \code{Hogwarts} $\to$ \code{Dumbledore}, \code{Malfoy}) \\
\bottomrule
\end{tabular}
\end{center}

\separator

\textbf{Analogical reasoning (parallelogram method):}

\begin{thmbox}[Parallelogram Model for Analogies]
To solve ``$a$ is to $b$ as $a^*$ is to ?'':
\[
\hat{b}^* = \arg\max_{x} \; \text{similarity}(x, \; a^* - a + b)
\]

\textbf{Examples:}
\begin{itemize}
  \item $\vec{\text{king}} - \vec{\text{man}} + \vec{\text{woman}} \approx \vec{\text{queen}}$
  \item $\vec{\text{Paris}} - \vec{\text{France}} + \vec{\text{Italy}} \approx \vec{\text{Rome}}$
\end{itemize}
\end{thmbox}

\begin{tipbox}[Caveats with Analogies]
\begin{itemize}
  \item Works mainly for \textbf{frequent words} and \textbf{certain relation types} (capitals, POS).
  \item Fails for rarer words, longer distances, and many semantic relations.
  \item Nearest neighbours may include query words themselves.
\end{itemize}
\end{tipbox}

\separator

\textbf{Bias in embeddings:}

\begin{tipbox}[Embeddings Reflect Corpus Biases]
\begin{itemize}
  \item ``father : doctor :: mother : ?'' $\to$ \code{nurse}
  \item ``man : computer programmer :: woman : ?'' $\to$ \code{homemaker}
\end{itemize}

\textbf{Implication:} Systems using biased embeddings (e.g., for hiring) can perpetuate discrimination.

\textbf{Historical embeddings} trained on different decades reveal changing cultural biases (e.g., gender stereotypes declining 1960--1990).
\end{tipbox}

\bigseparator

% --- Summary Table ---
\subsubsection{Summary: Sparse vs.\ Dense Methods}

\begin{center}
\renewcommand{\arraystretch}{1.3}
\begin{tabular}{@{} l p{5.5cm} p{5.5cm} @{}}
\toprule
 & \textbf{Count-Based (Sparse)} & \textbf{Word2Vec (Dense)} \\
\midrule
Representation & Co-occurrence counts (tf-idf, PMI) & Learned weight vectors \\
Dimensionality & $|V|$ & 50--1000 \\
Training & Direct matrix construction & SGD on prediction task \\
Synonymy & Separate dimensions & Similar vectors \\
Interpretability & High (dimensions = words) & Low (abstract dimensions) \\
Generalisation & Limited & Better (nearby = similar) \\
\bottomrule
\end{tabular}
\end{center}
